% jnb shared TikZ preamble.
% jnb/tikz.py substitutes the two marker lines below at compile time: the colour
% definitions are generated from assets/tokens.json, the font setup from tokens.json["fonts"]
% plus assets/fonts. The figure source is appended after \begin{document}.
% Compiled with lualatex --output-format=dvi, converted by dvisvgm --no-fonts.
\def\pgfsysdriver{pgfsys-dvisvgm.def}% pgf natively emits dvisvgm specials (no Ghostscript needed)
\documentclass[tikz,12pt,border=2pt]{standalone}
\makeatletter
\usepackage{amsmath}
\usepackage{fontspec}
\usepackage{unicode-math}
@@FONTS@@
\usepackage{xcolor}
@@COLORS@@
\usetikzlibrary{arrows.meta,calc,positioning,3d,perspective,fit}
\AtBeginDocument{\color{jnb-text-body}}

% ---- role macros (text and math): \inp{W_{t-1}}  \out{y}  \hl{0.16}
\newcommand{\inp}[1]{{\color{jnb-input}#1}}
\newcommand{\out}[1]{{\color{jnb-output}#1}}
\newcommand{\hl}[1]{{\color{jnb-highlight}#1}}
\newcommand{\third}[1]{{\color{jnb-third}#1}}
\newcommand{\fourth}[1]{{\color{jnb-fourth}#1}}
\newcommand{\negc}[1]{{\color{jnb-negative}#1}}

% ---- TikZ role styles
\tikzset{
  every picture/.append style={line cap=round, line join=round, >={Stealth[length=7pt,width=5pt]}},
  inp/.style={color=jnb-input},
  % `out` is also TikZ's "out=<angle>" key (to paths): keep that meaning when a value is given.
  out/.code={\if\relax\detokenize{#1}\relax\pgfkeysalso{color=jnb-output}\else\def\tikz@to@out{#1}\tikz@to@switch@on\fi},
  out/.default={},
  hl/.style={color=jnb-highlight},
  third/.style={color=jnb-third},
  fourth/.style={color=jnb-fourth},
  muted/.style={color=jnb-text-muted},
  inp arrow/.style={inp, line width=1pt, ->},
  out arrow/.style={out, line width=1pt, ->},
  leader/.style={hl, line width=1pt},
  jnb cell/.style={draw=jnb-grey-cell-edge, fill=jnb-grey-cell-fill, line width=0.8pt,
                   inner sep=0pt, outer sep=0pt, text=jnb-text-body, font=\small},
  jnb cell hl/.style={draw=jnb-highlight, fill=jnb-highlight!30!jnb-grey-cell-fill, line width=2pt},
  jnb cell inp/.style={draw=jnb-input, fill=jnb-input!22!jnb-surface-bg, line width=1.4pt, text=jnb-input},
  jnb cell out/.style={draw=jnb-output, fill=jnb-output!22!jnb-surface-bg, line width=1.4pt, text=jnb-output},
}

% =====================================================================
% \jnbcuboid{x}{y}{height}[options]
%   Isometric-style cuboid on the cell [x,x+1]x[y,y+1], drawn in the CURRENT 3D
%   coordinate system (set x={(..)}, y={(..)}, z={(..)} on a scope). Visible faces for a
%   viewer at large x, large y, large z: top, the x-facing face (left, greys.cuboid_left)
%   and the y-facing face (right, greys.cuboid_right).
%   options: gap (cell inset, 0.1)  base (z of the bottom, 0)  hl (flag: highlighted)
%            edge (colour name)  edge width  top/left/right (override face colours)
%   Highlight variant: highlight-coloured edges and translucent-looking highlight faces
%   (opacity .57 / .81 / .30 for top / left / right, pre-mixed with surface.bg so that
%   nothing shows through).
% =====================================================================
\newif\ifjnb@cuboidhl
\pgfkeys{
  /jnb/cuboid/.cd,
  gap/.initial=0.1, base/.initial=0, hl/.is if=jnb@cuboidhl,
  edge/.initial=jnb-grey-cell-edge, edge width/.initial=0.4pt,
  top/.initial=jnb-grey-cuboid-top, left/.initial=jnb-grey-cuboid-left, right/.initial=jnb-grey-cuboid-right,
}
\def\jnb@cuboidkey#1{\pgfkeysvalueof{/jnb/cuboid/#1}}
\def\jnb@cuboid#1#2#3#4{%
  \begingroup
  \jnb@cuboidhlfalse
  \pgfkeys{/jnb/cuboid/.cd,#4}%
  \pgfmathsetmacro\jnb@xa{#1+\jnb@cuboidkey{gap}}%
  \pgfmathsetmacro\jnb@xb{#1+1-\jnb@cuboidkey{gap}}%
  \pgfmathsetmacro\jnb@ya{#2+\jnb@cuboidkey{gap}}%
  \pgfmathsetmacro\jnb@yb{#2+1-\jnb@cuboidkey{gap}}%
  \pgfmathsetmacro\jnb@zz{\jnb@cuboidkey{base}}%
  \pgfmathsetmacro\jnb@zt{\jnb@cuboidkey{base}+#3}%
  \ifjnb@cuboidhl
    \def\jnb@ftop{jnb-highlight!57!jnb-surface-bg}%
    \def\jnb@fleft{jnb-highlight!81!jnb-surface-bg}%
    \def\jnb@fright{jnb-highlight!30!jnb-surface-bg}%
    \def\jnb@edge{jnb-highlight}\def\jnb@ew{1.4pt}%
  \else
    \edef\jnb@ftop{\jnb@cuboidkey{top}}\edef\jnb@fleft{\jnb@cuboidkey{left}}%
    \edef\jnb@fright{\jnb@cuboidkey{right}}%
    \edef\jnb@edge{\jnb@cuboidkey{edge}}\edef\jnb@ew{\jnb@cuboidkey{edge width}}%
  \fi
  % right face (y = yb), left face (x = xb), then top face
  \path[fill=\jnb@fright,draw=\jnb@edge,line width=\jnb@ew]
    (\jnb@xa,\jnb@yb,\jnb@zz)--(\jnb@xb,\jnb@yb,\jnb@zz)--(\jnb@xb,\jnb@yb,\jnb@zt)--(\jnb@xa,\jnb@yb,\jnb@zt)--cycle;
  \path[fill=\jnb@fleft,draw=\jnb@edge,line width=\jnb@ew]
    (\jnb@xb,\jnb@ya,\jnb@zz)--(\jnb@xb,\jnb@yb,\jnb@zz)--(\jnb@xb,\jnb@yb,\jnb@zt)--(\jnb@xb,\jnb@ya,\jnb@zt)--cycle;
  \path[fill=\jnb@ftop,draw=\jnb@edge,line width=\jnb@ew]
    (\jnb@xa,\jnb@ya,\jnb@zt)--(\jnb@xb,\jnb@ya,\jnb@zt)--(\jnb@xb,\jnb@yb,\jnb@zt)--(\jnb@xa,\jnb@yb,\jnb@zt)--cycle;
  \endgroup
}
\NewDocumentCommand{\jnbcuboid}{m m m O{}}{\jnb@cuboid{#1}{#2}{#3}{#4}}

% =====================================================================
% \jnbbars[options]{{row0 values},{row1 values},...}
%   Depth-sorted grid of cuboids (row index i -> x axis, column index j -> y axis).
%   Painter's order: increasing i+j, so nearer bars cover farther ones.
%   Height of a bar = min + scale * value.
%   options: scale (z units per unit value, 8)  min (flat-plate thickness, 0.2)
%            hl (list "i/j,i/j" of highlighted cells)  cuboid ({\jnbcuboid options})
%   Leaves \jnbbarsrows / \jnbbarscols defined; per-cell value is \jnbbarval{i}{j}.
% =====================================================================
\pgfkeys{
  /jnb/bars/.cd,
  scale/.initial=8, min/.initial=0.2, hl/.initial={}, cuboid/.initial={},
}
\def\jnbbarval#1#2{\csname jnb@v@#1@#2\endcsname}
\NewDocumentCommand{\jnbbars}{O{} m}{%
  \begingroup
  \pgfkeys{/jnb/bars/.cd,#1}%
  \edef\jnb@barscuboid{\pgfkeysvalueof{/jnb/bars/cuboid}}%
  \edef\jnb@hllist{\pgfkeysvalueof{/jnb/bars/hl}}%
  \def\jnb@nrows{0}\def\jnb@ncols{0}%
  \foreach \jnb@row [count=\jnb@ri from 0] in {#2}{%
    \global\let\jnb@nrows\jnb@ri
    \foreach \jnb@val [count=\jnb@ci from 0] in \jnb@row{%
      \expandafter\xdef\csname jnb@v@\jnb@ri @\jnb@ci\endcsname{\jnb@val}%
      \xdef\jnb@ncols{\jnb@ci}%
    }%
  }%
  \pgfmathtruncatemacro\jnb@nr{\jnb@nrows+1}%
  \pgfmathtruncatemacro\jnb@nc{\jnb@ncols+1}%
  \global\let\jnbbarsrows\jnb@nr \global\let\jnbbarscols\jnb@nc
  \ifx\jnb@hllist\@empty\else
    \foreach \jnb@h in \jnb@hllist{\expandafter\gdef\csname jnb@hl@\jnb@h\endcsname{1}}%
  \fi
  \pgfmathtruncatemacro\jnb@maxd{\jnb@nr+\jnb@nc-2}%
  \foreach \jnb@d in {0,...,\jnb@maxd}{%
    \foreach \jnb@i in {0,...,\jnb@nrows}{%
      \pgfmathtruncatemacro\jnb@j{\jnb@d-\jnb@i}%
      \ifnum\jnb@j>-1\relax\ifnum\jnb@j<\jnb@nc\relax
        \pgfmathsetmacro\jnb@hgt{\pgfkeysvalueof{/jnb/bars/min}+\pgfkeysvalueof{/jnb/bars/scale}*\jnbbarval{\jnb@i}{\jnb@j}}%
        \ifcsname jnb@hl@\jnb@i/\jnb@j\endcsname
          \edef\jnb@call{\noexpand\jnb@cuboid{\jnb@i}{\jnb@j}{\jnb@hgt}{hl,\jnb@barscuboid}}%
        \else
          \edef\jnb@call{\noexpand\jnb@cuboid{\jnb@i}{\jnb@j}{\jnb@hgt}{\jnb@barscuboid}}%
        \fi
        \jnb@call
      \fi\fi
    }%
  }%
  \endgroup
}

% =====================================================================
% \jnbmatrix[options]{{row0 cells},{row1 cells},...}
%   Grid of cells with the cell text; top-left corner of the matrix at the origin of the
%   current coordinate system, x to the right, rows going DOWN. Cell (i,j) is a rectangle
%   node named <name>-<i>-<j> (default name "M"), e.g. (M-3-4.east) for annotations.
%   options: name (M)  cw/ch (cell width/height, 1.7cm / 0.95cm)  hl (list "i/j,...",
%            highlighted cells, drawn last with the thick highlight border)
%   Leaves \jnbmatrixrows / \jnbmatrixcols defined.
% =====================================================================
\pgfkeys{
  /jnb/matrix/.cd,
  name/.initial=M, cw/.initial=1.7cm, ch/.initial=0.95cm, hl/.initial={},
}
\NewDocumentCommand{\jnbmatrix}{O{} m}{%
  \begingroup
  \pgfkeys{/jnb/matrix/.cd,#1}%
  \edef\jnb@mname{\pgfkeysvalueof{/jnb/matrix/name}}%
  \edef\jnb@mcw{\pgfkeysvalueof{/jnb/matrix/cw}}\edef\jnb@mch{\pgfkeysvalueof{/jnb/matrix/ch}}%
  \edef\jnb@hllist{\pgfkeysvalueof{/jnb/matrix/hl}}%
  \def\jnb@nrows{0}\def\jnb@ncols{0}%
  \foreach \jnb@row [count=\jnb@ri from 0] in {#2}{%
    \global\let\jnb@nrows\jnb@ri
    \foreach \jnb@val [count=\jnb@ci from 0] in \jnb@row{%
      \xdef\jnb@ncols{\jnb@ci}%
      \expandafter\xdef\csname jnb@mv@\jnb@ri @\jnb@ci\endcsname{\jnb@val}%
      \node[jnb cell,minimum width=\jnb@mcw,minimum height=\jnb@mch]
        (\jnb@mname-\jnb@ri-\jnb@ci) at ({(\jnb@ci+.5)*\jnb@mcw},{-(\jnb@ri+.5)*\jnb@mch}) {\jnb@val};
    }%
  }%
  \pgfmathtruncatemacro\jnb@nr{\jnb@nrows+1}\pgfmathtruncatemacro\jnb@nc{\jnb@ncols+1}%
  \global\let\jnbmatrixrows\jnb@nr \global\let\jnbmatrixcols\jnb@nc
  \ifx\jnb@hllist\@empty\else
    \foreach \jnb@h in \jnb@hllist{\expandafter\jnb@mhl\jnb@h\relax}%
  \fi
  \endgroup
}
\def\jnb@mhl#1/#2\relax{%
  \node[jnb cell hl,minimum width=\jnb@mcw,minimum height=\jnb@mch]
    at (\jnb@mname-#1-#2.center) {\csname jnb@mv@#1@#2\endcsname};%
}
